\section{Theorie}
\subsection{Elektronenhülle und Atomspektren}
Die potenzielle Energie eines Elektrons in der Atomhülle nimmt mit zunehmendem Abstand zum Atomkern zu, welcher durch die Hauptquantenzahl~$n$ beschrieben wird. Elektronen können durch die Absorption von Licht Energie aufnehmen und zwischen diesen diskreten Zuständen wechseln. Diese Übergänge erzeugen die charakteristischen Absorptions- und Emissionsspektren, da jedes Element nur spezifische Wellenlängen absorbieren oder emittieren kann.
Die Energie $E$ elektromagnetischer Strahlung ist gegeben durch:
\begin{equation}
  E = h f = \frac{h c}{\lambda}
  \label{eq:energie}
\end{equation}
wobei $h = \SI{6,62607015e-34}{J \cdot s}$ das Plancksche Wirkungsquantum, $f$ die Frequenz, $\lambda$ die Wellenlänge und $c = \SI{299792458}{m/s}$ die Vakuumlichtgeschwindigkeit ist \cite{nist_constants}. Wie in Gleichung~\ref{eq:energie} zu erkennen, sind Energie und Wellenlänge invers proportional zueinander.

\subsection{Balmer-Serie}
Das Wasserstoffatom emittiert Licht nur bei bestimmten diskreten Wellenlängen, die als Emissionsspektrallinien beobachtet werden. Diese lassen sich in Serien einteilen, die jeweils durch einen festen unteren Quantenzustand charakterisiert sind. Die Frequenzen der emittierten Strahlung werden durch die Rydberg-Formel beschrieben \cite{demtroeder_experimentalphysik3}:
\begin{equation}
  f_{n_2 \rightarrow n_1} = R_{\ce{H}} \cdot c \cdot \left(\frac{1}{n_1^2} - \frac{1}{n_2^2}\right)
  \label{eq:balmer-frequenz}
\end{equation}
Hierbei sind $n_1$ und $n_2$ die Hauptquantenzahlen des unteren bzw. oberen Zustands. Für die Balmer-Serie gilt $n_1 = 2$ und $n_2 = 3,4,5,\dots$. Die Konstante $R_{\ce{H}} = \SI{1,09737e7}{m^{-1}}$ ist die Rydberg-Konstante \cite{nist_constants}. Die einzelnen Linien werden als $H_\alpha$, $H_\beta$, $H_\gamma$ bezeichnet und liegen im sichtbaren Spektralbereich (ca.~\SIrange{400}{800}{nm}).

\subsection{Interferenz am optischen Gitter}
Trifft Licht auf ein optisches Gitter mit Gitterkonstante~$g$, entstehen Interferenzmaxima der Ordnung~$k$ unter einem Winkel~$\alpha$ gegenüber der optischen Achse. Mit dem Abstand~$l$ zwischen Gitter und Schirm sowie dem Abstand~$x$ vom zentralen Maximum gilt \cite{hecht_optik}:
\begin{align}
  \sin(\alpha) &= \frac{k \cdot \lambda}{g}, \quad \tan(\alpha) = \frac{x}{l} \notag \\
  \implies \lambda &= \frac{\sin\!\left(\arctan\!\left(\dfrac{x}{l}\right)\right) \cdot g}{k}
  \label{eq:optisches-gitter}
\end{align}

\subsection{Grundlagen der optischen Absorptionsspektroskopie}
Die optische Absorptionsspektroskopie untersucht die wellenlängenabhängige Schwächung von Licht beim Durchgang durch eine Probe. Wird monochromatisches Licht der Ausgangsintensität $I_0(\lambda)$ durch eine Küvette der Schichtdicke~$d$ geleitet, so wird es auf die Intensität $I_1(\lambda)$ abgeschwächt. Die Abschwächung hängt von der Konzentration $c_0$ und dem molaren Extinktionskoeffizienten $\epsilon(\lambda)$ ab. Dieser Zusammenhang wird durch das Lambert-Beer'sche Gesetz beschrieben \cite{skript}:
\begin{equation}
  I_1(\lambda) = I_0(\lambda) \cdot e^{-\epsilon(\lambda)\cdot c_0 \cdot d}
  \label{eq:lambert-beer}
\end{equation}
Im sichtbaren und UV-Bereich werden dabei vor allem Elektronen in diskrete Energieniveaus angeregt.

\subsection{Das Teilchen im eindimensionalen Kasten}
Beim Modell des Teilchens im eindimensionalen Kasten wird ein Elektron in einem Potentialkasten der Länge~$L$ mit unendlich hohen Wänden betrachtet. Durch Lösung der stationären Schrödinger-Gleichung ergeben sich quantisierte Energiezustände \cite{atkins_physical_chemistry}:
\begin{equation}
  E_n = \frac{n^2 h^2}{8 m_e L^2}, \quad n \in \mathbb{N}
  \label{eq:kasten-energie}
\end{equation}
wobei $m_e$ die Elektronenmasse und $n$ die Hauptquantenzahl ist. Für ein System mit $N$~$\pi$-Elektronen sind aufgrund des Pauli-Prinzips im Grundzustand die niedrigsten $N/2$ Niveaus doppelt besetzt. Der HOMO-LUMO-Übergang hat die Energiedifferenz:
\begin{equation}
  \Delta E = \frac{h^2}{8 m_e L^2}(N+1)
  \label{eq:energie-potentialtopf}
\end{equation}
Gleichsetzen mit $\Delta E = hc/\lambda$ (Gl.~\ref{eq:energie}) liefert die Kastenlänge:
\begin{equation}
  L = \sqrt{\frac{h \cdot \lambda_{\text{max}} \cdot (N+1)}{8 m_e c}}
  \label{eq:kastenlänge}
\end{equation}
